\section{The Constitution of a Traditional Society: The Capo}

\paragraph{The Supreme Leader}

\begin{quotex}
This is the concluding section of \textbf{Guido De Giorgio}`s description of the establishement of a traditional society. Here, he develops the idea of the supreme leader, above the three castes, whose primary duty is the maintenance of Tradition.

This section, as a whole, needs to be compared to the similar efforts by Guenon, Coomaraswamy, and Evola to describe the respective roles of Spiritual Authority and Temporal Power. However, De Giorgio's work comes from a completely different place, from a deep understanding of the principles involved. The other three try to make an ``argument", by referring to various scriptures or historical situations. Guido De Giorgio speaks from the heart, from what he knows; he can be judged on that basis. That is how Tradition is to be recovered. 

Interspersed within the main topic, as always there are various spiritual gems to be recovered. We can call attention in particular to the notion that a man creates his own life. Also, pace those who want to create a superman, De Giorgio believes that man qua man must vanish in order for something higher to appear. 

\end{quotex}
We could repeat at this point the exhortation that Homer put into the mouth of wise Ulysses when the Greeks appeared to abandon the Trojan adventure and that even Aristotle cited at the end of Book XII of the Metaphysics: ``let there be one ruler." There is first of all the great law of analogy according to which the only leader in time must correspond to the supreme unity in eternity and since God is pure contemplation and only the cognitive fruition of aseity can be conceived in Him, so the Leader inversely will make a pure activity of his life dedicated to the maintenance of his command on the earth.

He holds temporal power, exercises it uncontested, and is the supreme authority whose domain embraces all the active life in its multiple aspects and over which he is the regulator and judge. His function is therefore integrative and his work is strictly dependent on Justice, the highest of the virtues.

If on the one hand, while being the supreme leader, he is above all the prince of the Warriors, his power is exercised over all indiscriminately in the temporal environment because he could not go beyond it without failing in his task and his purpose.

Although he arose from the Warrior caste, in order to carry out his task he is above the three castes and so must, while remaining strictly within the scope of the active life, serve as the harmonious base of the development of the first two castes, from the first receiving consecration, and serving for the protection and the maintenance of the traditional order of the second. In order to preserve inviolate his sovereignty, he will have to guard the Temple that is the foundation of every tradition and defend the Priests who nourish with their spiritual influence the body of the truths accessible only to the contemplatives. While depending spiritually on them, he is absolutely autonomous in his domain and the Priests in this owe him obedience as all men belonging to the other two castes, nor can or should the spiritual insert itself into the temporal and compromise the pure sphere of contemplation which does not have to undergo profanation of that sort.

The unity of command represents the centralizing equilibrium of the active life that loses in the Leader its particularistic and widespread character in order to rise to a level analogous to the Edenic state where everything is put before man who is truly the king of creation. Therefore the Leader vanishes only before God where his sovereignty is annulled in the Lordship of all the lords: everything in this life, he will be in the other that which he will have merited according to the meaning it will have had from God. ``Caesar I was, Justinian I am…" [Dante, Paradise vii]. In the face of God the greatest victory will be that which he will have been able to bring back on himself and on what he will have accomplished on earth, nothing will matter for him except the works of justice, rectification, equilibrium, and peace. Since the Leader must always seek peace for his subjects in analogy with that same thing that in the contemplative order is the foundation of every realization and all aspire to that type of activity that develops them. The maintenance of this equilibrium is reserved to the Warriors who will always make war an instrument of peace and not of pure and simple conquest. The unity of the empire avoids the unproductive diffusions of the active life concentrating its energies toward a single end where all disparities and partial imbalances are evened out, therefore if in the Warriors, as was said, inner asceticism is the necessary base for the development of their activity, this is even more strict in the Leader who must annul in himself any individuality whatsoever in order to make justice triumph and the rule secure the empire.

He knows that he is the supreme ruler of the active life and that the active vibration emanates from him that, in analogy with the divine, is extended to all the branches of humanity of which he constitutes the propelling and purifying center through which the backward wave [onda di ritorno] obtains new energy for new possibilities of development in a perpetually mobile and fecund circle of flux and reflux.

We are purely in the sphere of action in which all men necessarily live and to which none are exempt except those who can and must dedicate themselves to the contemplative life: what each man accomplishes in it, and the way in which he accomplishes it, has a very great importance for the equilibrium of human perfection that would be maimed and amputated if an element is removed from that current of active homogeneity that constitutes temporal modality.

All traditions in fact reserve the contemplative life to those for whom it is appropriate and who are worthy of it, while they condemn asceticism and separation from the world for those who lie outside it obeying a sense of inadmissible egoism: those are the weak and the lifeless and they avoid the responsibility of worldly life through real insufficiency. It will be well to insist on that: everyone brings and creates, so to speak, his environment and by this word we do not mean only the closest circle, but the entire temporal extent of human life with everything which is positive and negative in it for the individual. Making use here of an image, we could say that every man has a passage of existence that is not only his world, but all the world that is for him what he is for the world itself: in this passage that starts from the center of a circumference and is contained between two divergent rays which, having as base a segment of the circumference itself, form a cone, there is everything that man wants and that he does not want, his fortune, his misfortune, his battle, his dispute, his victory or his downfall. The essential thing for him is to rise up in ascending planes from the circumference to the center, from the base of the cone to the vertex that represents the greatest of the minimum, i.e., God, while reassimilating all multiplicity gradually into a point only in which he is annulled, because he, dispelling ignorance, will recognize himself in everything, even and above all, in that to which he is most opposite and will recompose in a perfect equilibrium all the antagonisms up to the point of extinguishing them in the highest vortex of his asceticism, or better said, of his recovery, which is God.

When we employ this term, we mean everything that is truly of God in God, truly and not fallaciously, i.e., all the superhuman states of union with God that man can realize only starting from himself in order to reach that which is no longer himself, man before, no longer man afterwards.

The terrestrial level is one experience and not another, that is necessary to start from a base of disablement in order to reach an apex of absolute immutability: these are the two greatest terms in the sphere of that which appears in the seat of human relativity: a fallacious reality, man, a genuine reality, God, the former only illusory, the latter, supreme certainty. But placing oneself at an integrative point of view, one is reality, God, who is no longer a goal which one reaches except through a dissipation of ignorance, as whoever believes real the phantasms of the fog that are non-existent provided the sun dispels and dissolves it.

\paragraph{The Role of the Leader}

\begin{quotex}
In this second of three parts, \textbf{Guido De Giorgio} expands on the role of the Leader in a traditional society. It is instructive to compare De Giorgio's point of view with Evola's idea of the state; De Giorgio includes the spiritual or contemplative aspect which brings in some subtle changes. Nevertheless, Di Giorgio also recognizes that the active life can also be a path to realization.

The emphasis on the law is also important, such Guenon pointed out that the western medieval tradition did not have a sacred law. There was, and is, canon law, but that is at a lesser level. What De Giorgio is describing is more akin to sharia law in Islam, or the complex rituals and requirements required of Hindu castes, and even the ancient Romans whose daily lives were regulated by rites and taboos. Abandoning such laws, in his view, does not lead to more freedom, but rather to is opposite: it leads to servitude to one's passions and lower nature. 

\end{quotex}
On the other hand, by not knowing himself man does not know God because only the knowledge of his illusory existence permits him to realize the essential truth of God in a process that moves from negation to affirmation, from human nothingness to divine wholeness, not from a fragment to a totality, but from a non-totality to totality, from illusory human unity to real divine unity, not from duality to unity, but from non-duality to unity. This knowledge entails an effective realization that in its turn annuls the earthly sphere, the environment in which man lives, putting all its elements on a plane of absolute equivalence, i.e., neutralizing, dissipating the sense of otherness that derives from the possessive conceit concluding in the ontological absurdity ``I am I".

The process of realization instead denies the first person, affirms the second, unifies the first and third, and definitively puts itself beyond this union, in what we could call the Fourth Sublime, the Fourth Absolute.

This is the ideal plan of transhumanization presented as a fieri while it is an esse, and it is the decisive model of everything in principle and therefore also and above all of the Leader who is not such if he does not take on the experience of his subjects in the purity of the distributive standard and in the unification of sovereign power. Therefore his function is, so to say, the restoration of hierarchy, the consecration of the human body whose members are the subjects while he is the vivifying center that accomodates the innumerable expressions of the active life, purifies them in the renovating power of traditional principles whose deposit the Priests hold and returns them almost equipped with a more intense rhythm and clad with the united seal.

In this sense he the first and the last of his subjects because his subordination to divine principles is dependent on his capacity to centralize the forms of active life while expending his personality in a decisive totalization. The relationship between the Leader and his subjects consists in this, that he is in all and all are in him with an absolute reciprocity that makes of the entire active life a veritable sacrifice, i.e., an imperfection made sacred through the gift that is made of it to God in the person of the Leader. He accepts and gives, and what he offers is always himself: he is responsible for how much his subjects accomplish because in him everything merges in order to receive the supreme consecration. The asceticism of the Leader is in a sense superior to that of the Warriors because it is vaster and more integral, accomodating in himself all the developments of the active life in order to maintain them within the traditional ambit with the spiritual authority of the Priests and the power of the Warriors, from the former, the conscience, and from the latter, the personality so that ascetically the last in front of God is the first in front of men. His activity, obeying the law of justice, is thus purged of any personal motives whatsoever because, by becoming Leader, he consecrated himself to the coming of peace on earth and governs in order to make it respected and he provides the harmony of traditional society with a constant effort of rectification. As God is the ruler of all the worlds, he is the ruler of the earth and of the men who inhabit it, the supreme regulator of their activity without passing beyond his domain which is that of the active life, without which he could be in open conflict with his function and purpose as temporal leader occupying the contemplative sphere that is to the Priests reserved. Harmony must precisely consist in the separation of the two powers that, on earth, and on earth alone, where activity is in force, conserving two distinct spheres while it no longer exists in an ultra-worldly stage that is the starting point for the integral realization of the divine.

The subjects owe the Leader obedience and respect; obedience because he directs human affairs, analogously with God who rules the creative universality, and respect because his function is the application of justice, i.e., the rectification of the disparities produced by the multiformity of men's actions and protects them and guides them to the realization of happiness in whose terminal and plenary stage is Edenic perfection.

Human activity when brought back to the purpose of rite is purified from every contaminating egoism and the man who has fulfilled his temporal stage in prefect adherence to the norm of justice, regains the absolute preeminence over the other beings of whom he becomes again lord and he enjoys the Great Peace in the exercise of his liberty released from every abuse. It is the Leader's precise function to tend to this reintegration of the divine state, directing all the forms of earthly activity toward a progressive purgation and acting so that man can fulfill what is granted to him on earth without actions limiting and precluding contemplation but, regulated and ordered, more than a snag and a chain, is itself a mode of liberation. And he can become it only if conducted according to justice which is the greatest of the virtues, especially for the Leader who must prefer it more than any other virtue because his action neutralizes the will leading every form of activity back into its channel, preventing abuses of power and giving true liberty to the subjects, i.e., the spiritual practice that is gained with obedience to the law.

Here man is liberated by obeying and not otherwise: whoever does not understand that will always be subjugated. In fact, in order to obey it is necessary to understand the value of the law that is placed to redeem men from the appetites that enslave and to make him able to exercise his freedom effectively. The Leader must be vigilant in respect to the law in order to facilitate the gaining of freedom and he himself will be liberated to the extent he maintains the limits of his temporal function without ever passing beyond them. The law in itself does not know how to nor can it make itself respected, but it needs whoever establishes it, and since it is a pure norm higher than men, it must be imposed: in this sense strength accompanies justice, nor can it be separated and the Leader must make the law respected with strength. But from submission to the law, freedom is born that, being purely of the spiritual order and belonging to the contemplative life, arises only from the perfection of this, like a crowning achievement, an apical stage that is the triumph of man over himself.

Man calls himself free only when he again becomes the son of truth and does not recognize subjection other than this to the point of confusing himself, identifying with a further state which will be that of pure unity. But in the domain of the active life, it is not possible to be liberated other than by obeying the law personified in the Leader who must make it respected with strength, not through an affirmation of supremacy— which is to be excluded in a traditional society—but to prevent the subjects from falling into servitude by breaking the law. The temporal power is inclined to the achievement of freedom and is established for this, so that, by imposing the law, the immoderate will is avoid, the bestial part in man is restrained, and it is started precisely to be free. To regulate the active life can signify nothing other that is positively fertile than to permit contemplative practice in the measure conceded to each one by his spirit. In civil life, freedom is therefore obedience and in the contemplative life it is liberation: up to this point it is understood that man is not so much born free as he becomes free, and better said, born free, he can maintain his freedom or fall into servitude: most men, by forgetting freedom, which is a divine gift, prefer servitude to their passions, and this they call freedom and are opposed to the monarchical regime because they see in the Leader a man like them and not a model who, in order to make himself respected, must necessarily personify himself: one is not liberated, but becomes liberated, and naturally he becomes it because he is it, but ignorance prevents him from knowing why man is born free: only for the fulfillment of his earthly experience and not for the exaltation of that which is least human in him. He was free at his birth as the son of God, but, straying in order to follow his impulse, which he fallaciously calls freedom, he has become again a slave: in order to become free again, i.e., a son of God, he must put himself under the regulative norm of the active life and to regain the freedom of the contemplative life.

\paragraph{Conclusion}

\begin{quotex}
This is the conclusion to the role of the Leader in a traditional society. It needs to be read meditatively, and not with the intent to refute or falsify it. If read in silence, without the infusion of personal opinions, something very different may be heard. 

\end{quotex}
We said that man first denies himself and then affirms himself: he denies himself in the law and affirms himself in love, which is only the supplementary knowledge of the truth. Now the active life cannot be abandoned to the disordered will of the appetites which would make it an end in itself, while it is only a means whose purpose is clearly determined by the temporal power embodied in the Leader. In order to be truly such, it is necessary that this power be employed by a single man so that every subject contemplates himself in his sovereign, his better part, to whom he must submit himself, the reason that regulates every activity restraining them within their limits and preventing them from gravely obstructing the attainment of the pure values of the spirit.

If the natural end of man is heaven, the starting point is the earth, so that the necessity for making its foundation organized by a liberating action that acts especially on the obstructing element, the body, which abandoned to itself, precludes every possibility of overcoming and attainment. That which is corporeal remains subordinated to the psychic element that is transparently vague through its elusiveness and diffusivity: the virtues inherent in the active life, that the Leader rules and directs, have as their field of action the psychic complex that must be contained and limited by reason as if by a supreme command. But reason is in a certain sense two-edged to the extent that on the one hand it dominates and restrains the lower faculties and on the other, it inserts into a higher order of reality where the intellect extends and, while being human, is the mediator between the human and the divine: likewise the Leader rules in the sign of justice so that a higher stage becomes a reality, whose domain belongs more specially to the Priests, but it is necessary for the Warriors to facilitate this perfection.

Given that, it indicates that war is the condition of peace as long as it is truly conducted by those for whom it is their only given mission and within the limits that make it justifiable and necessary.

In the earthly and human domain, law must make use of force and the symbol and the means of justice is two-edged sword, one of which strikes and points toward the lower, which the other points higher and is bloodless, the point representing the peak of unification of the higher and the lower, the supreme acme, death, or resolution Thus what seemingly appears cruel, is not always so in reality, and the justice that imposes itself with force is often the necessary vestibule of the highest sphere of love. The activity of the Leader in human and earthly life is analogous to that of God, whose direct instrument he is, in order to coordinate the modes of activity pursuant to justice, truth, and love. No one more than the Leader knows that by punishing others he punishes himself and that with an act of justice he fulfills an achievement of love and truth: if he does not know this, if the ignores his purpose, he is not truly a leader, but a slave of his own cupidity, since he cannot command if he is not aware of the raison d'.tre of his authority, of the divine character of his mission.

A Warrior among the Warriors he must be an ascetic among the ascetics: so he will reach the point to rule others by how much more he will rule himself and to defeat his enemies by defeating himself, and to fight the small wars by fighting the Great War, and to achieve external peace by how much more profoundly and perfectly he will have achieved interior peace, knowing that earthly life is only a stage, but the most important of all, of which he is the orderer. His responsibility is very great before men but especially before God since by betraying his creatures he betrays his Creator and by neglecting justice in order to follow his own impulse he contradicts the nobility of the position that he occupies and abuses the power that is given to him: then he will truly be last among men and servant among the servants. Since there is no justice without a judge, so the Leader is also judge among men and here his task is still more serious and difficult: he must be very attentive and not exceed the scope in which his activity unfolds, since as judge he can affect the case, but not the knowledge of those who have accomplished it. He knows in fact that everything that happens may not happen and that there is nothing open to censure that others accomplish that he himself cannot accomplish, since men are the mirror of man and he, as Leader, is the mirror of the time in which he lives and is also the supreme responsible man for it.

His more or less great adherence to tradition is measured by the degree of greater or lesser spirituality of the era in which he lives and of which he is the judge. He must dematerialize the world, bring active life back within its normal limits, to make action a means and not an end, and contemplation an end and not a means, and force the pure expression of law and power, the assertion of his humility before God, and justice the implacable rectification, and the gift of himself the constant offering.

He is in the world to redeem it, to purify it in the name of God, for the love of God, because he, with the Warriors, is the protector of the Temple and his mission is the respect of the sacred law to which he is subject and through which he is the Leader, since everything that comes from God must return to God, and what is accomplished on earth is only for the dignity of the true man who is a son of God. His true glory, his true conquest will have facilitated it, prepared it, realized the coming of the kingdom of God, having transformed life in a rite redeeming man from his state of servitude, leading him to liberty through to liberation, making of the active life a connecting bridge to the beatitude of contemplation where the divine is only realized, ransoming not oppressing, spiritualizing not materializing.

Only in that way will he be the God's chosen and his power will receive that consecration from which all authority comes. He knows that man and the world are of God and that everything must flow into the great traditional axis that is sustained by sacred knowledge of which he is the most pious and aware supporter. The active life that he rules and normalizes must be contained within the limits that permit each man to rise to the sphere of contemplation where true life begins, the divine life, and his mission unfolds in the temporal sphere for eternity, on earth for heaven, the center of the innumerable efforts that flows from every point into a single wave, in a single rhythm, for the glory of the one before Whom everything is nothing and the nothing is all.

Supreme bundler of fasces [Fascificatore], the Leader redeems the transience of man and the world with the vivifying rhythm of tradition of which he must be the most intransigent defender in order to establish the constitutive unity of castes, the harmony of principles, the respect of standards, the convergence of all the efforts in the fulfillment of his very high mission, so that the world is truly the mark, the crown, and the sign of God.



\flrightit{Posted on 2013-07-21 by Cologero }

\begin{center}* * *\end{center}

\begin{footnotesize}\begin{sffamily}



\texttt{Jason-Adam on 2013-07-21 at 12:26 said: }

In the Christian tradition, the capo existed in the form of the Byzantine Emperor, and later the Russian Tsar who was the absolute ruler over both spiritual and temporal realms……..in the East the Emperor protected defended the church but the church was still subject to his laws. In the West the Capo never emerged due to conflicts between the Popes and the Emperors……


\hfill

\texttt{Cologero on 2013-07-22 at 13:07 said: }

J-A, do you suppose that Guido De Giorgio is proposing a thought experiment along the lines of Plato's Republic? Although the domain of the Capo is activity, unlike God, he can never be pure act, so there is always non-being or privation associated with his rule. The stronger Leaders can come closer to the ideal.


\hfill

\texttt{JA on 2013-07-23 at 12:16 said: }

Of course de Giorgio is discussing an ideal, but it is our task to make the ideal real….if not now then when the time is right……..

It is also a good question whether or not the Capo could ever actually exist, perhaps Justinian or Charlemagne came close ?


\hfill

\texttt{JA on 2013-07-29 at 10:42 said: }

From my first readings de Giorgio's ideas read very much like that of the Byzantine-Russian practice of symphonia where church and state were separate but linked in that the Tsar was anointed by the patriarch and expected to be a good Christian at the risk of losing his legitimacy from God, but while subject to the church in spiritual matters the Emperor was ruler of the church in temporal realms. 

In the western history I'm trying to understand if Roman Catholic doctrine would allow for a Capo or would the Pope claim too much power in temporal affairs that an absolute monarch would never be able to exist ?


\hfill

\texttt{Cologero on 2013-07-29 at 12:39 said: }

Tomberg addresses this very question, as it is enshrined in the arcana of the Pope and the Emperor.

For example:

``The authors of the Middle Ages could not imagine Christianity without an Emperor."


\hfill

\texttt{JA on 2013-07-29 at 14:17 said: }

Tomberg follows the Catholic conception but is his vision of the Emperor the same as de Giorgio's ? I'll have to read this essay several more times before I am able to conclude for sure but now it's seeming to me that the Capo has more power than the Holy Roman Emperor (for example the priests are subject to him in the temporal sphere, whereas the Church claimed that a King or and Emperor had no power over a clergyman, the clergy answer to the Pope alone, hence no state could be absolute because the clergy constituted a caste that was not in service to the monarch) – later on the arrangement changed more in the Kings favour (Spanish Baroque period, France under Louis XIV) but that was the decline of Catholicism already…..

basically my worry about the Catholic system of government is that the duality of heads is unstable (Pope vs Emperor) and allows for rebellions of kings (as what happened with protestantism). The Orthodox model seems more stable but then its problem is that it become too national (greek church, russian church all under their own Tsars) and fails to create an Imperium as Rome was meant to do…


\hfill

\texttt{Logres on 2013-07-30 at 10:27 said: }

The Orthodox apparently have a tradition that St. Feodor of Tomsk was Alexander I, who staged his own assassination.

\url{http://www.johnsanidopoulos.com/2012/11/was-saint-feodor-of-tomsk-really.html}


\hfill

\texttt{Logres on 2013-07-30 at 10:53 said: }

Would de Giorgio's doctrine imply that, even in our time, such a person necessarily exists, however hidden or concealed? It seems to me that it would – God would be sustaining his link to the world through the sacred body of the King. King Alfred seems like a close approximation to the ideal, in the West.


\hfill

\texttt{scardanelli on 2013-07-30 at 11:57 said: }

Logres,

You jogged my memory, and I recalled that I read of this same concept in the introduction to Sedir's Initiations:

``There is a Prince of this Earth –delegate of Lucifer; and there is also the Lord of this Earth –delegate of Christ…[T]he Lord of the Earth is the one among the servants, soldiers or friends of

Christ who possesses the qualities and the right quantity of Lights most suited to link our terrestrial life to the Word directly. He receives the orders of Christ, executes them, then transmits the hopeful aspirations of the earth's inhabitants to Him."


\hfill

\texttt{Ash on 2013-08-02 at 01:58 said: }

An essential question to ask here must be, who judges the Capo? Or rather, can anyone decide when he has failed in his role? Presumably it would be the duty of the Contemplative caste to make these judgements and of the Active (warrior) caste to carry them out and replace the Leader. We might think of what happened during the Empire when Emperors such as Caligula, Nero, and Elagabulus were killed by the Praetorian guard (admittedly, money and preservation of power had more to do with those than concern for Traditional principles). A troublesome situation arises when the Capo begins to feud with the other two higher castes, such as Henry VIII and the Catholic Church. While it is obviously not de Giorgio's principle, I must admit to some unease as much of the language used here could easily be hijacked to justify totalitarian and decidedly anti-Traditional forms of ideology and tyranny. There are many out there eager to grab the reigns of power and ignore the voices of those who remind them that this power is meant to be exercised within its proper dominion and no further.


\hfill

\texttt{Cologero on 2013-08-03 at 14:01 said: }

Perhaps, Ash, tomorrow's conclusion will clarify some things. However, it seems you are thinking in terms of the modern mind and the examples you give are hardly traditional. It's true, there is literally no one to judge the Capo since he is God's representative and no one can judge God. Of course, the modern mind regards that as a trick, some kind of ruse to keep the people under control, the opiate of the masses. The modern mind denies the very possibility of God manifesting through some particular man, whereas the traditional mind can see things no other way: the cosmic order is in continuity and harmony with the social order.

Jung made an interesting point, for example, that the Pharaoh was the only individuated person and the mass was undifferentiated. The latter state would have been more like that described by Julian Jaynes and thoughts would have seemed to them to originate from the outside, from ``gods". Hence, the command of the Pharaoh would have had the same force. We may suppose that the consciousnesses of the scribes and priest were similar to the Pharaoh's to some extent. Jung and Serrano claim that the ``persona" is more marked in the West and that may account for the turmoil.

In the Ancient City, the leader was a priest-king with divine authority. Unfortunately, we don't have sufficient documentation, but Fustel de Coulanges claims that at some point the function of the Chief Priest and King got separated as the Warrior caste asserted more control. The relevant period for Rome was the Kingdom. For example, the second king Numa also led the Priests. It is interesting that the way the King was selected is nearly identical to the way Popes are selected today. Again, that documentation was lost, to the extent that in the Republic and the Empire, the ancient laws became incomprehensible.

I guess the point is that the traditional system is not simply one alternative among many, as some sort of political party among many others. It would require an intellectual conversion, a real change in consciousness.


\hfill

\texttt{Ash on 2013-08-03 at 18:51 said: }

Cologero, would you not say there is a concern to address in terms of the abuse of such authority? Clearly the boundaries of the Capo are defined and as such the truly differentiated man would be following them, but we also have accounts of corrupt Emperors and Pharaohs. I suppose in some ways this question is very cart-before-the-horse, since presumably a traditional force in society would establish itself before the question of a Capo ever arose. But if, as de Giorgio says, the priests consecrate the Capo even as he comes from the warriors, that would seem to mean that they can prevent unworthy individuals from attaining such position, and would it not simply be a further fulfillment of this role that they act as a check on the Capo (if ideally an unnecessary one). To put it more plainly, I think we are all generally suspicious of individuals claiming that they are God's representative and that they must be followed as His own voice, precisely because they are generally scammers and anti-Traditional. We, as students of Tradition, must strive more than anyone to be sure that those who guide and teach are authentic and acting in accordance with their duties and the Truth. 

On the point of the King I can provide some assistance, as I did some in depth studying on the religious heritage of Rome for university papers last year. My main secondary source was a book called Religions of Rome (Beard, North, and Price). It makes the case that the political Rex and the Rex Sacrocum were always separate positions in the regal era. During the early Republic and onward, the Rex Sacrocum was a member of the Pontifical college and as such himself under the authority of the Pontifex Maximus. This is however a new development, as it was previously assumed by many that the Rex Sacrocum at some point simply took over the religious duties of the King. Your example of Numa is ideal though: assuming separate functions, Numa would seem to fulfill the role of Capo and thus the Rex Sacrocum would be subordinate to him, being the head priest.


\hfill

\texttt{Cologero on 2013-08-03 at 19:34 said: }

If the Capo has the qualities that De Giorgio describes, then there is no abuse. The Roman king selection followed an interesting process. The Senate elected the King from among them, so presumably they should know him well. He was presented to the people who could refuse him. Finally, the priest had to wait for the omen to indicate the new king had divine authority. So the king required approval from the three castes. Is that not sufficient for you? In your studies on the Roman Kingdom, why did it eventually fail and get replaced by the Republic?


\hfill

\texttt{Ash on 2013-08-03 at 21:12 said: }

That would seem more than sufficient for choosing the ruler, yes. I am thinking more about degeneration over time (as the castes themselves face as well). Perhaps it would be useful to reflect on past incidences of this degeneration. Keeping Rome as our example, the last King was remembered as a tyrant who forever swore the city off that system of government. 

As to why the Kingdom fell…the simple answer is that historians are left in many ways to guesswork. We know of the story of the rape of Lucretia and the fall of Tarquinius Superbus, but I would personally think that the ``other" myth of Romulus' death may provide insight. In this account, Romulus is murdered by the Patres, the body which became the Senate. They then dismember his body and attempt to hide the murder by claiming that he became a god, who was worshiped as Quirinus. We know that the Kings were generally chosen from outside the Curae, the collections of families from the Roman population. Thus the six legendary Kings consist of Latins, Sabines, and Etruscans, none of whom directly succeed fathers. This suggests that conflict always existed at some level between the Patres/Senators and the Kings who were meant to deliberate among them as outsiders, as would be expected. As an aside, we also know that the senatorial and priestly bodies were overlapping to a great extent, and that the major priests would often also be major political figures, particularly in the later Republic. 

Having derived these points, and knowing that the Republic was established following a rebellion of the Senatorial class against a King who had become considered tyrannical, we may conclude that this conflict came to a head. History being written by the winners, we may assume that the Senatorial patrician classes decided that they were better off with the familial competition than regal moderation. If we are to examine this with de Giorgio's vocabulary, this means that the warriors (the Senators and Equestrians alone had the property and wealth to be able to train and arm themselves to protect Rome) overthrew the Capo who had been previously approved by them. We do not know who the ``people" (the plebs and plebeians) supported in this struggle, though later generations remembered the overthrow as a great event and shared the Senate's hatred of Kings. Perhaps the previously mentioned defeat of the rebellion of the Kshatriyas in the Hindu epics is an example of a similar conflict which was won by the other side.


\hfill

\texttt{Saladin on 2013-08-03 at 22:28 said: }

Ash@ ``We do not know who the ``people" (the plebs and plebeians) supported in this struggle"

Ash, I am sure you meant Plebs and Proletariat, right?

As you said, it's not too much of a stretch to think that ``Senatorial patrician classes decided that they were better off with the familial competition than regal moderation" and they replaced a degenerate king with a Republic. As for the masses, I don't think their taking side would have made that much of a difference at that period of time. Even though ``bread and circus" was invented for them much later, during the Empirical period, a different version of ``bread and circus" was enough to keep them preoccupied during those early Republican period.


\hfill

\texttt{Ash on 2013-08-04 at 00:24 said: }

Saladin, I did mean pleb and plebeian. The two words are not wholly interchangeable, as pleb and proletarian would probably be. Pleb means a labourer and man without property who is working in labour or perhaps a trade, whereas plebeian simply denotes one who is not part of the original patrician families. The patrician/plebeian distinction actually lost importance rather early during the Republic, when the ``nobiles"/new man distinction between those who had held consulships and those who had not gained importance. Many major figures during the Republic and almost all of those in the Empire were plebeian, but not plebs as they had attained high standing. Augustus' familial links with the Julians aside, there were no real patrician Emperors. It's just a distinction between terms and sometimes gets confused with catchalls like nobility or aristocracy, and for here I wouldn't worry too much about it.


\hfill

\texttt{Avery Morrow on 2013-08-04 at 05:09 said: }

Ash, Fustel de Coulanges distinguishes the concept of rex from tyrannus, a king who rules without religious affirmation. All Emperors were tyranni because their authority came from the secularized SPQR and not the patrician faith. (They weren't categorized that way by their contemporaries but they would have been recognized as such in the preclassical city — a great demonstration of how much Rome changed over the centuries.) Regardless of whether rex and rex sacrocum were different positions, it should go without saying that they both had the symbolic quality of rex.

A comparison to the Far East is interesting.


\end{sffamily}\end{footnotesize}
